\chapter{告子章句上}
\kaishi*

\jiazhu[1]{告子者，告姓也。子，男子之通稱也，名不害。兼治儒墨之道者，嘗學於孟子，而不能純徹性命之理。《論語》曰：「子罕言命。」謂性命之難言也。以告子能執弟子之問，故以題篇。}

\section{告子論性猶杞柳}

告子曰：「性猶杞柳也，義猶桮棬也。以人性爲仁義，猶以杞柳爲桮棬。」\jiazhu[1]{告子以爲人性爲才幹，義爲成器，猶以杞柳之木爲桮棬也。杞柳，柜柳也。一曰杞，木名也。《詩》云：「北山有杞。」桮棬，桮素也。}

孟子曰：「子能順杞柳之性而以爲桮棬乎？將戕賊杞柳而後以爲桮棬也？」\jiazhu[1]{戕，猶殘也。《春秋傳》曰：「戕舟發梁。」子能順完杞柳不傷其性而成桮棬乎？將以斤斧殘賊之，乃可以爲桮棬乎？言必殘賊也。}

「如將戕賊杞柳而以爲桮棬，則亦將戕賊人以爲仁義與？」\jiazhu[1]{孟子言以人身爲仁義，豈可復殘傷其形體乃成仁義邪？明不可比桮棬也。}

「率天下之人而禍仁義者，必子之言夫！」\jiazhu[1]{以告子轉性以爲仁義，若轉木以成器，必殘賊之。故言率人以禍仁義者，必子之言夫！歎辭也。}\par \jiazhu[1]{章指言：養性長義，順夫自然；殘木爲器，變而後成。告子道偏，見有不純；仁內義外，違人之端。孟子拂之，不假以言也。}

\section{告子論性猶湍水}

告子曰：「性猶湍水也，決諸東方則東流，決諸西方則西流。人性之無分於善不善也，猶水之無分於東西也。」\jiazhu[1]{湍者，圜也，謂湍湍瀠水也。告子以喻人性若是水也，善惡隨物而化，無本善不善之性也。}

孟子曰：「水信無分於東西，無分於上下乎？人性之善也，猶水之就下也。人無有不善，水無有不下。今夫水，搏而躍之，可使過顙；激而行之，可使在山。是豈水之性哉？其勢則然也。人之可使爲不善，其性亦猶是也。」\jiazhu[1]{孟子曰：水誠無分於東西，故決之而往也。水豈無分於上下乎？水性但欲下耳。人性生而有善，猶水欲下也。所以知人皆有善性，似水無有不下者也。躍，跳；顙，額也。人以手跳水，可使過顙；激之，可令上山。皆迫於勢耳，非水之性也。人之可使爲不善，非順其性也，亦妄爲利欲之勢所誘迫耳。猶是水也，言其本性非不善也。}\par \jiazhu[1]{章指言：人之欲善，猶水好下。迫勢激躍，失其素眞。是以守正性者爲君子，隨曲拂者爲小人也。}

\section{告子論生之謂性}

告子曰：「生之謂性。」\jiazhu[1]{凡物生同類者，皆同性。}

孟子曰：「生之謂性也，猶白之謂白與？」\jiazhu[1]{猶見白物皆謂之同白，無異性也。}

曰：「然。」\jiazhu[1]{告子曰然。}

「白羽之白也，猶白雪之白；白雪之白，猶白玉之白與？」\jiazhu[1]{孟子以爲羽性輕，雪性消，玉性堅。雖俱白，其性不同。問告子：子以三白之性同邪？}

曰：「然。」\jiazhu[1]{告子曰：然，誠以爲同也。}

「然則犬之性猶牛之性，牛之性猶人之性與？」\jiazhu[1]{孟子言：犬之性豈與牛同所欲？牛之性豈與人同所欲乎？}\par \jiazhu[1]{章指言：物雖有性，性各殊異。惟人之性，與善俱生。赤子入井，以發其誠。告子一之，知其麤矣。孟子精之，是在其中。}

\section{告子論仁內義外}

告子曰：「食色，性也。仁，內也，非外也；義，外也，非內也。」\jiazhu[1]{人之甘食悅色者，人之性也。仁由內出，義在外也，不從己身出也。}

孟子曰：「何以謂仁內義外也？」\jiazhu[1]{孟子怪告子是言也。}

曰：「彼長而我長之，非有長於我也。猶彼白而我白之，從其白於外也，故謂之外也。」\jiazhu[1]{告子言：見彼人年長大，故我長敬之。長大者非在於我也，猶白色見於外也。}

曰：「異於白馬之白也，無以異於白人之白也。不識長馬之長也，無以異於長人之長與？且謂長者義乎？長之者義乎？」\jiazhu[1]{孟子曰：長異於白。白馬、白人同謂之白，可也。不知敬老馬無異於敬老人邪？且謂老者爲有義乎？將謂敬老者爲有義乎？敬老者己也，何以爲外也？}

曰：「吾弟則愛之，秦人之弟則不愛也，是以我爲悅者也，故謂之內。長楚人之長，亦長吾之長，是以長爲悅者也，故謂之外也。」\jiazhu[1]{告子曰：愛從己，則己心悅，故謂之內；所悅喜老在外，故曰外。}

曰：「耆秦人之炙，無以異於耆吾炙。夫物則亦有然者也。然則耆炙亦有外與？」\jiazhu[1]{孟子曰：耆炙同等，情出於中。敬楚人之老與敬己之老亦同，己情往敬之。雖非己炙，同美，故曰物則有然者也。如耆炙之意，豈在外邪？言楚、秦喻遠也。}\par \jiazhu[1]{章指言：事雖在外，行其事者皆發於中。明仁義由內，所以曉告子之惑也。}

\section{孟季子問義內}

孟季子問公都子曰：「何以謂義內也？」\jiazhu[1]{季子亦以爲義外也。}

曰：「行吾敬，故謂之內也。」\jiazhu[1]{公都子曰：以敬在心而行之，故言內。}

「鄉人長於伯兄一歲，則誰敬？」\jiazhu[1]{季子曰：敬誰也？}

曰：「敬兄。」\jiazhu[1]{公都子曰：當敬兄也。}

「酌則誰先？」\jiazhu[1]{季子曰：酌酒則先酌誰？}

曰：「先酌鄉人。」\jiazhu[1]{公都子曰：當先鄉人。}

「所敬在此，所長在彼，果在外，非由內也。」\jiazhu[1]{季子曰：所敬者兄也，所酌者鄉人也。如此，義果在外，不由內也。果猶竟也。}

公都子不能答，以告孟子。\jiazhu[1]{公都子無以答季子之問。}

孟子曰：「敬叔父乎？敬弟乎？彼將曰：『敬叔父。』曰：『弟爲尸，則誰敬？』彼將曰：『敬弟。』子曰：『惡在其敬叔父也？』彼將曰：『在位故也。』子亦曰：『在位故也。』庸敬在兄，斯須之敬在鄉人。」\jiazhu[1]{孟子使公都子答季子如此。言弟以在尸位，故敬之；鄉人在賓位，故先酌之耳。庸，常也。常敬在兄，斯須之敬在鄉人也。}

季子聞之，曰：「敬叔父則敬，敬弟則敬，果在外，非由內也。」\jiazhu[1]{隨敬所在而敬之，果在外。}

公都子曰：「冬日則飲湯，夏日則飲水。然則飲食亦在外也？」\jiazhu[1]{湯水雖異名，其得寒温者中心也。雖隨敬之所在，亦中心敬之，猶飲食從人所欲，豈可復謂之外也？}\par \jiazhu[1]{章指言：凡人隨形，不本其原；賢者達情，知所當然。季子信之，猶若告子；公都受命，然後乃理。}

\section{公都子論性之善惡}

公都子曰：「告子曰：『性無善無不善也。』」\jiazhu[1]{公都子道告子以爲人性在化，無本善不善也。}

「或曰：『性可以爲善，可以爲不善。是故文、武興，則民好善；幽、厲興，則民好暴。』」\jiazhu[1]{公都子曰：或人以爲可教以善不善，亦由告子之意也。故文、武聖化之起，民皆喜爲善；幽、厲虐政之起，民皆好暴亂。}

「或曰：『有性善，有性不善。是故以堯爲君而有象；以瞽瞍爲父而有舜；以紂爲兄之子，且以爲君，而有微子啓、王子比干。』」\jiazhu[1]{公都子曰：或人者以爲人各有性善惡，不可化移。堯爲君，象爲臣，不能使之爲善；瞽瞍爲父，不能化舜爲惡；紂爲君，又與微子、比干有兄弟之親，亦不能使此二子爲不仁。是亦各有性也。}

「今曰『性善』，然則彼皆非與？」\jiazhu[1]{公都子曰：告子之徒，其論如此。今孟子曰人性盡善，然則彼之所言皆非邪？}

孟子曰：「乃若其情，則可以爲善矣，乃所謂善也。若夫爲不善，非才之罪也。」\jiazhu[1]{若，順也。性與情相爲表裏，性善勝情，情則從之。《孝經》曰：「此哀戚之情。」情從性也。能順此情，使之善者，眞所謂善也。若隨人而強作善者，非善者之善也。若爲不善者，非所受天才之罪，物動之故也。}

「惻隱之心，人皆有之；羞惡之心，人皆有之；恭敬之心，人皆有之；是非之心，人皆有之。惻隱之心，仁也；羞惡之心，義也；恭敬之心，禮也；是非之心，智也。仁義禮智，非由外鑠我也，我固有之也，弗思耳矣。故曰：『求則得之，舍則失之。』或相倍蓰而無算者，不能盡其才者也。」\jiazhu[1]{仁義禮智，人皆有其端，懷之於內，非從外消鑠我也。求存之則可得而用之，舍縱之則亡失之矣。故人之善惡，或相倍蓰，或至於無算者，不得相與計多少，言其絕遠也。所以惡乃至是者，不能自盡其才性也。故使有惡人，非天獨與此人惡性。其有下愚不移者，譬如被疾不成之人，所謂童昏也。}

「《詩》曰：『天生蒸民，有物有則。民之秉彝，好是懿德。』孔子曰：『爲此詩者，其知道乎！故有物必有則；民之秉彝也，故好是懿德。』」\jiazhu[1]{《詩·大雅·烝民》之篇。言天生衆民，有物則有所法則，人法天也。民之秉彝，彝，常也，常好美德。孔子謂之知道，故曰人皆有善也。}\par \jiazhu[1]{章指言：天之生人，皆有善性，引而趨之，善惡異衢，高下相懸，賢愚舛殊。尋其本者，乃能一諸。}

\section{孟子論富歲凶歲}

孟子曰：「富歲，子弟多賴；凶歲，子弟多暴。非天之降才爾殊也，其所以陷溺其心者然也。」\jiazhu[1]{富歲，豐年也；凶歲，飢饉也。子弟，凡人之子弟也。賴，善；暴，惡也。非天降下才性與之異也，以飢寒之厄陷溺其心，使爲惡者也。}

「今夫麰麥，播種而耰之，其地同，樹之時又同，浡然而生，至於日至之時，皆熟矣。雖有不同，則地有肥磽，雨露之養，人事之不齊也。」\jiazhu[1]{麰麥，大麥也。《詩》云：「貽我來麰。」言人性之同如此麰麥，其不同者，人事雨澤有不足，地之有肥磽耳。磽，薄也。}

「故凡同類者，舉相似也，何獨至於人而疑之？聖人與我同類者。」\jiazhu[1]{聖人亦人也，其相覺者以心知耳。蓋體類與人同，故舉相似也。}

「故龍子曰：『不知足而爲屨，我知其不爲蕢也。』屨之相似，天下之足同也。」\jiazhu[1]{龍子，古賢者也。雖不知足小大作，屨者猶不更作蕢。蕢，草器也。以屨相似，天下之足略同故也。}

「口之於味，有同耆也。易牙先得我口之所耆者也。如使口之於味也，其性與人殊，若犬馬之與我不同類也，則天下何耆皆從易牙之於味也？至於味，天下期於易牙，是天下之口相似也。」\jiazhu[1]{人口之所耆者相似，故皆以易牙爲知味，言口之同也。}

「惟耳亦然。至於聲，天下期於師曠，是天下之耳相似也。」\jiazhu[1]{耳亦猶口也。天下皆以師曠爲知聲之微妙也。}

「惟目亦然。至於子都，天下莫不知其姣也。不知子都之姣者，無目者也。」\jiazhu[1]{目亦猶耳也。子都，古之姣好者也。《詩》云：「不見子都，乃見狂且。」儻無目者，乃不知子都好耳。言目之同耳。}

「故曰：口之於味也，有同耆焉；耳之於聲也，有同聽焉；目之於色也，有同美焉。至於心，獨無所同然乎？」\jiazhu[1]{言人之心性皆同也。}

「心之所同然者何也？謂理也，義也。聖人先得我心之所同然耳。故理義之悅我心，猶芻豢之悅我口。」\jiazhu[1]{心所同耆者，義理也。理者，得道之理。聖人先得理義之要耳。理義之悅心，如芻豢之悅口，誰不同也？草牲曰芻，穀養曰豢。}\par \jiazhu[1]{章指言：人稟性俱有好憎，耳目口心所悅者同。或爲君子，或爲小人，猶麰麥不齊，雨露使然也。孟子言是，所以勗而進之。}

\section{孟子論牛山之木}

孟子曰：「牛山之木嘗美矣，以其郊於大國也，斧斤伐之，可以爲美乎？是其日夜之所息，雨露之所潤，非無萌櫱之生焉，牛羊又從而牧之，是以若彼濯濯也。人見其濯濯也，以爲未嘗有材焉，此豈山之性也哉？」\jiazhu[1]{牛山，齊之東南山也。邑外謂之郊。息，長也。濯濯，無草木之貌。牛山未嘗盛美，以在國郊，斧斤牛羊使之不得有草木耳，非山之性無草木也。}

「雖存乎人者，豈無仁義之心哉？其所以放其良心者，亦猶斧斤之於木也，旦旦而伐之，可以爲美乎？其日夜之所息，平旦之氣，其好惡與人相近也者幾希。」\jiazhu[1]{存在也。言雖在人之性，亦猶此山之有草木也。人豈無仁義之心邪？其日夜之思欲息長仁義，平旦之志氣，其好惡凡人皆有與賢人相近之心。幾，豈也。豈希，言不遠也。}

「則其旦晝之所爲，有梏亡之矣。梏之反覆，則其夜氣不足以存；夜氣不足以存，則其違禽獸不遠矣。人見其禽獸也，而以爲未嘗有才焉者，是豈人之情也哉？」\jiazhu[1]{旦晝，晝日也。其所爲萬事，有梏亂之，使亡失其日夜之所息也。梏之反覆，利害干其心，其夜氣不能復存也。人見惡人禽獸之行，以爲未嘗有善才性，此非人之情也。}

「故苟得其養，無物不長；苟失其養，無物不消。孔子曰：『操則存，舍則亡；出入無時，莫知其鄉。』惟心之謂與？」\jiazhu[1]{誠得其養，若雨露於草木，法度於仁義，何有不長？誠失其養，若斧斤牛羊之消草木，利欲之消仁義，何有不盡也？孔子曰：持之則在，縱之則亡，莫知其鄉。鄉猶里，以喻居也。獨心爲若是也。}\par \jiazhu[1]{章指言：秉心持正，使邪不干，猶止斧斤不伐牛山，山則木茂，人則稱仁也。}

\section{孟子論王之不智}

孟子曰：「無或乎王之不智也。」\jiazhu[1]{王，齊王也。或，怪也。時人有怪王不智而孟子不輔之，故言此也。}

「雖有天下易生之物也，一日暴之，十日寒之，未有能生者也。吾見亦罕矣，吾退而寒之者至矣，吾如有萌焉何哉？」\jiazhu[1]{種易生之草木五穀，一日暴溫之，十日陰寒以殺之，物何能生？我亦希見於王。既見而退，寒之者至，謂左右佞諂順意者多。譬諸萬物，何由得有萌芽生也？}

「今夫弈之爲數，小數也；不專心致志，則不得也。」\jiazhu[1]{弈，博也，或曰圍棋。《論語》曰：「不有博弈者乎？」數，技也。雖小技，不專心則不得也。}

「弈秋，通國之善弈者也。使弈秋誨二人弈，其一人專心致志，惟弈秋之爲聽；一人雖聽之，一心以爲有鴻鵠將至，思援弓繳而射之。雖與之俱學，弗若之矣。爲是其智弗若與？曰：非然也。」\jiazhu[1]{有人名秋，通一國皆謂之善弈，曰弈秋。使教二人弈，其一人惟秋所善而聽之，其一人念欲射鴻鵠，故不如也。爲是謂其智不如也？曰：非也，以不致志也。故齊王之不智，亦若是。}\par \jiazhu[1]{章指言：弈爲小數，不精不能。一人善之，十人惡之，雖竭其道，何由智哉？《詩》云：「濟濟多士，文王以寧。」此之謂也。}

\section{孟子論舍生取義}

孟子曰：「魚，我所欲也；熊掌，亦我所欲也。二者不可得兼，舍魚而取熊掌者也。生亦我所欲也，義亦我所欲也。二者不可得兼，舍生而取義者也。」\jiazhu[1]{熊掌，熊蹯也，以喻義；魚以喻生也。}

「生亦我所欲，所欲有甚於生者，故不爲苟得也；死亦我所惡，所惡有甚於死者，故患有所不辟也。」\jiazhu[1]{有甚於生者，謂義也。義者不可苟得。有甚於死者，謂無義也，不苟辟患也。莫甚於生，則苟利而求生矣；莫甚於死，則可辟患不擇善，何不爲耳？}

「由是則生而有不用也，由是則可以辟患而有不爲也。是故所欲有甚於生者，所惡有甚於死者。非獨賢者有是心也，人皆有之，賢者能勿喪耳。」\jiazhu[1]{有不用，不用苟生也；有不爲，不爲苟惡而辟患也。有甚於生，義甚於生也；有甚於死，惡甚於死也。凡人皆有是心，賢者能勿喪亡之也。}

「一簞食，一豆羹，得之則生，弗得則死。嘑爾而與之，行道之人弗受；蹴爾而與之，乞人不屑也。」\jiazhu[1]{人之餓者得此一器食可以生，不得則死。嘑爾，猶呼爾，咄啐之貌也。行道之人，道中凡人，以其賤己，故不肯受也。蹴，蹋也，以足踐蹋與之。乞人不絜之，亦由其小，故輕而不受也。}

「萬鍾則不辯禮義而受之，萬鍾於我何加焉？爲宮室之美、妻妾之奉、所識窮乏者得我與？」\jiazhu[1]{言一簞食則貴禮，至於萬鍾則不復辯別有禮義與不。鍾，量器也。萬鍾於己身何加益哉？己身不能獨食萬鍾也，豈不爲廣美宮室、供奉妻妾、施與所知之人窮乏者？}

「鄉爲身死而不受，今爲宮室之美爲之；鄉爲身死而不受，今爲妻妾之奉爲之；鄉爲身死而不受，今爲所識窮乏者得我而爲之，是亦不可以已乎？此之謂失其本心。」\jiazhu[1]{鄉者不得簞食而食則身死，尚不受也。今爲此三者爲之，是不亦可以止乎？所謂失其本心也。}\par \jiazhu[1]{章指言：舍生取義，義之大者也。簞食萬鍾，用有輕重。縱彼納此，蓋違其本。凡人皆然，君子則否，所以殊也。}

\section{孟子論求放心}

孟子曰：「仁，人心也；義，人路也。舍其路而弗由，放其心而不知求，哀哉！」\jiazhu[1]{不行仁義者，不由路，不求心者也，可哀憫哉！}

「人有雞犬放，則知求之；有放心而不知求。學問之道無他，求其放心而已矣。」\jiazhu[1]{人知求雞狗，莫知求其心者，惑也。學問所以求之。}\par \jiazhu[1]{章指言：由路求心，爲得其本；追逐雞狗，務其末也。學以求之，詳矣。}

\section{孟子論不知類}

孟子曰：「今有無名之指屈而不信，非疾痛害事也。如有能信之者，則不遠秦楚之路，爲指之不若人也。」\jiazhu[1]{無名之指，手之第四指也。蓋以其餘指皆有名，無名指者，非手之用指也。雖不疾痛妨害於事，猶欲信之，不遠秦楚，爲指不若人故也。}

「指不若人，則知惡之；心不若人，則不知惡。此之謂不知類也。」\jiazhu[1]{心不若人，可惡之大者也，而反惡指，故曰不知其類也。類，事也。}\par \jiazhu[1]{章指言：舍大惡小，不知其要。憂指忘心，不鄉於道。是以君子惡之也。}

\section{孟子論養身}

孟子曰：「拱把之桐梓，人苟欲生之，皆知所以養之者。至於身，而不知所以養之者，豈愛身不若桐梓哉？弗思甚也。」\jiazhu[1]{拱，合兩手也；把，以一手把之也。桐梓，皆木名也。人皆知灌漑而養之，至於養身之道當以仁義，而不知用，豈於身不若桐梓哉？不思之甚也。}\par \jiazhu[1]{章指言：莫知養身而養樹木，失事違務，不得所急，所以誡未達者也。}

\section{孟子論養其大者}

孟子曰：「人之於身也，兼所愛。兼所愛，則兼所養也。無尺寸之膚不愛焉，則無尺寸之膚不養也。」\jiazhu[1]{人之所愛則養之，於身也，一尺一寸之膚養相及也。}

「所以考其善不善者，豈有他哉？於己取之而已矣。」\jiazhu[1]{考知其善否，皆在己之所養也。}

「體有貴賤，有小大。無以小害大，無以賤害貴。養其小者爲小人，養其大者爲大人。」\jiazhu[1]{養小則害大，養賤則害貴。小，口腹也；大，心志也。頭頸，貴者也；指拇，賤者也。不可舍貴養賤也。務口腹者爲小人，治心志者爲大人。}

「今有場師，舍其梧檟，養其樲棘，則爲賤場師焉。」\jiazhu[1]{場師，治場圃者。場以治穀，圃，園也。梧、桐、檟、梓，皆木名。樲棘，小棘，所謂酸棗也。言此以喻人舍大養小，故曰賤場師也。}

「養其一指而失其肩背，而不知也，則爲狼疾人也。」\jiazhu[1]{謂醫養人疾，治其一指而不知其肩背之有疾，以至於害之，此爲狼籍亂不知治疾之人也。}

「飲食之人，則人賤之矣，爲其養小以失大也。飲食之人無有失也，則口腹豈適爲尺寸之膚哉？」\jiazhu[1]{飲食之人，人所以賤之者，爲其養口腹而失道德耳。如使不失道德，存仁義以往，不嫌於養口腹也。故曰：口腹豈但爲肥長尺寸之膚邪？亦爲懷道者也。}\par \jiazhu[1]{章指言：養其行，治其正，俱用智力，善惡相厲。是以君子居處思義，飲食思禮也。}

\section{公都子問大人小人}

公都子問曰：「鈞是人也，或爲大人，或爲小人，何也？」\jiazhu[1]{鈞，同也。言有大有小，何也？}

孟子曰：「從其大體爲大人，從其小體爲小人。」\jiazhu[1]{大體，心思禮義；小體，縱恣情欲。}

曰：「鈞是人也，或從其大體，或從其小體，何也？」\jiazhu[1]{公都子言：人何獨有從小體也？}

曰：「耳目之官不思，而蔽於物。物交物，則引之而已矣。心之官則思，思則得之，不思則不得也。此天之所與我者。先立乎其大者，則其小者弗能奪也。此爲大人而已矣。」\jiazhu[1]{孟子曰：人有耳目之官不思，故爲物所蔽。官，精神所在也，謂人有五官六腑。物，事也，利欲之事來交引其精神。心官不思善，故失其道而陷爲小人也。此乃天所與人情性。先立乎其大者，謂生而有善性也。小者，情欲也。善勝惡，則惡不能奪。}\par \jiazhu[1]{章指言：天與人性，先立其大，心官思之，邪不乖越，故謂之大人也。}

\section{孟子論天爵人爵}

孟子曰：「有天爵者，有人爵者。仁義忠信，樂善不倦，此天爵也；公卿大夫，此人爵也。」\jiazhu[1]{天爵以德，人爵以祿。}

「古之人修其天爵，而人爵從之。今之人修其天爵，以要人爵；既得人爵，而棄其天爵，則惑之甚者也。」\jiazhu[1]{人爵從之，人爵自至也。以要人爵，要求也。得人爵，棄天爵，惑之甚也。}

「終亦必亡而已矣。」\jiazhu[1]{棄善忘德，終必亡之。}\par \jiazhu[1]{章指言：古修天爵，自樂之也；今要人爵，以誘時也。得人棄天，道之忌也。惑以招亡，小人事也。}

\section{孟子論良貴}

孟子曰：「欲貴者，人之同心也。人人有貴於己者，弗思耳。人之所貴者，非良貴也。趙孟之所貴，趙孟能賤之。」\jiazhu[1]{人皆同欲貴之心，人人自有貴者在己身，不思之耳。在己者，謂仁義廣譽也。凡人之所貴富，故曰非良貴者。趙孟，晉卿之貴者，能貴人，又能賤人。人所自有者，他人不能賤之也。}

「《詩》云：『既醉以酒，既飽以德。』言飽乎仁義也，所以不願人之膏粱之味也；令聞廣譽施於身，所以不願人之文繡也。」\jiazhu[1]{《詩·大雅·既醉》之篇。言飽德者，飽仁義之於身，身之貴者也，不願人膏粱矣。膏粱，細粱如膏者也。文繡，繡衣服也。}\par \jiazhu[1]{章指言：所貴在身，人不知求。膏粱文繡，己之所優。趙孟所貴，何能比之？是以君子貧而樂也。}

\section{孟子論仁勝不仁}

孟子曰：「仁之勝不仁也，猶水勝火。今之爲仁者，猶以一杯水救一車薪之火也，不熄，則謂之水不勝火。此又與於不仁之甚者也，亦終必亡而已矣。」\jiazhu[1]{水勝火，取水足以制火。一杯水何能勝一車薪之火也？以此謂水不勝火。爲仁者亦若是，則與作不仁之甚者也。亡猶無也，亦終必無仁矣。}\par \jiazhu[1]{章指言：爲仁不至，不反諸己，謂水勝火，熄而後已。不仁之甚，終必亡矣。爲道不卒，無益於賢也。}

\section{孟子論仁在乎熟之}

孟子曰：「五穀者，種之美者也。苟爲不熟，不如荑稗。夫仁，亦在乎熟之而已矣。」\jiazhu[1]{熟，成也。五穀雖美，種之不成，不如荑稗之草，其實可食。爲仁不成，猶是也。}\par \jiazhu[1]{章指言：功毀幾成，人在慎終。五穀不熟，荑稗是勝。是以爲仁，必其成也。}

\section{孟子論教射與學道}

孟子曰：「羿之教人射，必志於彀；學者亦必志於彀。」\jiazhu[1]{羿，古之工射者。彀，張也，弩向包的者，用思要時也。學者志道，猶射者之張也。}

「大匠誨人，必以規矩；學者亦必以規矩。」\jiazhu[1]{大匠，攻木之工。規，所以爲圓也；矩，所以爲方也。誨，教也。教人必須規矩，學者以仁義爲法式，亦猶大匠以規矩者也。}\par \jiazhu[1]{章指言：事各有本，道有所隆。彀張規矩，以喻爲仁。學不爲仁，猶是二教，失其法而行之也。}
